% CVPR 2026 Paper Template
\documentclass[10pt,twocolumn,letterpaper]{article}

\usepackage{cvpr}
\usepackage{times}
\usepackage{epsfig}
\usepackage{graphicx}
\usepackage{amsmath}
\usepackage{amssymb}
\usepackage{booktabs}
\usepackage{multirow}
\usepackage{xcolor}
\usepackage{hyperref}

% Paper ID (assigned during submission)
\def\cvprPaperID{****}
\def\confName{CVPR}
\def\confYear{2026}

\begin{document}

\title{Topology-Aware Crack and Spalling Segmentation in Concrete Dams\\via Dynamic Snake Convolution and Skeleton Recall Loss}

\author{
Anonymous Author(s)\\
{\tt\small anonymous@institution.edu}
}

\maketitle

\begin{abstract}
% Abstract goes here (< 200 words)
\end{abstract}

\section{Introduction}
% Introduction content

\section{Related Work}
\subsection{Concrete Defect Segmentation}
\subsection{Topology-Aware Segmentation}
\subsection{Dynamic Snake Convolution}

\section{Method}
\subsection{Overview}
\subsection{DSCFormer Architecture}
\subsubsection{SegFormer-B2 Backbone}
\subsubsection{Dynamic Snake Convolution Branch}
\subsection{Skeleton Recall Loss}
\subsection{Composite Loss Function}

\section{Experiments}
\subsection{Dataset: DamSegment}
\subsection{Implementation Details}
\subsection{Main Comparison}
\subsection{Per-Tier Analysis}
\subsection{Ablation Study}
\subsection{Qualitative Results}

\section{Discussion}

\section{Conclusion}

{\small
\bibliographystyle{ieee_fullname}
\bibliography{refs}
}

\end{document}
